\chapter{Implementación}
\label{chap:implementacion}

\section{Análisis de tecnologías}
\label{sec:analisis_tecnologias}

La implementación del sistema requiere tecnologías compatibles con la
arquitectura existente de UVLHub y con el ecosistema Flamapy. Su selección se ha
realizado teniendo en cuenta la integración con la plataforma, la ejecución del
generador en los entornos necesarios y la disponibilidad de un entorno
reproducible para el desarrollo y el despliegue.

\subsection*{Necesidad técnica 1: Integración con la plataforma existente}

El generador debe integrarse en UVLHub, una plataforma web desarrollada con
Flask~\cite{flask}. La solución debe aprovechar sus rutas, plantillas Jinja2 y
mecanismos de sesión, sin introducir cambios estructurales innecesarios en la
aplicación.

\subsection*{Necesidad técnica 2: Desarrollo de la herramienta generadora}

La lógica de generación debe seguir la arquitectura de extensiones de
Flamapy. Para ello, se necesita un componente independiente que encapsule la
configuración del generador y la operación encargada de producir modelos de
características. La conversión posterior a UVL se realiza mediante
\texttt{UVLWriter}.

\subsection*{Necesidad técnica 3: Entornos de ejecución}

La generación ordinaria se ejecuta en el navegador mediante Pyodide y
WebAssembly. Cuando se solicita la comprobación de satisfacibilidad, la
generación y los reintentos se ejecutan en el servidor, que dispone de las
dependencias necesarias para utilizar PySAT. Ambos recorridos emplean la misma
operación \texttt{GenerateFeatureModel}.

\subsection*{Necesidad técnica 4: Desarrollo y despliegue reproducibles}

El sistema debe disponer de un entorno reproducible para el desarrollo local y
de una versión accesible durante las revisiones del proyecto. Esto requiere
contenerizar UVLHub y gestionar las dependencias de los componentes que se
ejecutan tanto en el navegador como en el servidor.

\subsubsection*{Comparativa de tecnologías del backend}

\begin{itemize}
    \item \textbf{Flask:} \textit{Framework} utilizado por UVLHub sobre el que se
    ha integrado el generador. Permite aprovechar directamente sus rutas,
    plantillas y mecanismos de sesión.

    \item \textbf{Django:} \textit{Framework} web con numerosas funcionalidades
    integradas. Su utilización habría requerido adaptar la arquitectura
    existente de UVLHub para incorporar el generador.

    \item \textbf{FastAPI:} \textit{Framework} orientado al desarrollo de APIs.
    Su adopción para esta funcionalidad habría añadido otro entorno web a una
    aplicación que ya utiliza Flask.
\end{itemize}

\subsubsection*{Comparativa de tecnologías del front-end}

\begin{itemize}
    \item \textbf{Jinja2, JavaScript y Bootstrap~\cite{bootstrap}:} Son las
    tecnologías utilizadas por UVLHub para construir las vistas y gestionar
    la interacción. Permiten integrar los pasos del asistente en la interfaz
    existente.

    \item \textbf{React:} Biblioteca de interfaces basada en componentes.
    Incorporarla habría exigido una arquitectura adicional para las pantallas
    del generador y su comunicación con UVLHub.

    \item \textbf{Vue.js:} \textit{Framework} de interfaces reactivas. No se
    seleccionó porque los pasos del asistente ya se integran en las plantillas
    y los recursos JavaScript de UVLHub.
\end{itemize}

\subsubsection*{Análisis de la base de datos}

El generador integrado no requiere una base de datos propia. Durante el
asistente, la configuración se conserva temporalmente en la sesión de Flask;
los modelos obtenidos se ofrecen para su visualización o descarga. No se ha
añadido persistencia para guardar estas configuraciones ni los resultados de
la generación. Esta decisión no modifica los mecanismos de almacenamiento
utilizados por otras funcionalidades de UVLHub.

\subsubsection*{Plataformas de despliegue}

\begin{itemize}
    \item \textbf{Render~\cite{render}:} Plataforma utilizada para ofrecer una
    versión accesible de UVLHub durante las revisiones. Su elección responde a
    la facilidad de configuración y a la experiencia previa del autor.

    \item \textbf{Servicios \textit{cloud} tradicionales (AWS, Azure y Google
    Cloud):} Ofrecen más posibilidades de configuración de infraestructura,
    que no eran necesarias para el despliegue de revisión previsto.

    \item \textbf{Heroku:} Habría permitido un despliegue similar, pero se
    disponía de experiencia previa y de una configuración establecida con
    Render.
\end{itemize}

\subsubsection*{Tecnologías finales seleccionadas}

\begin{itemize}
    \item \textbf{Servidor y lógica de generación:} Python, Flask, Flamapy y
    PySAT para la comprobación de satisfacibilidad solicitada por el usuario.

    \item \textbf{Interfaz:} Plantillas Jinja2, JavaScript y Bootstrap,
    siguiendo la estructura existente de UVLHub.

    \item \textbf{Ejecución ordinaria en el navegador:} Pyodide y WebAssembly,
    junto con las dependencias empaquetadas para el generador.

    \item \textbf{Persistencia propia del generador:} No se incorpora
    almacenamiento permanente de configuraciones ni de modelos generados.

    \item \textbf{Entorno local y despliegue de revisión:} Docker y Docker
    Compose para el desarrollo local, y Render para ofrecer la aplicación a los
    tutores durante las revisiones.
\end{itemize}

\section{Tecnologías utilizadas}
\label{sec:tecnologias}

Esta sección describe cómo se han utilizado las tecnologías seleccionadas y
en qué componentes de la aplicación aparecen.

\subsection*{Backend y lógica del servidor}

El backend de UVLHub está desarrollado en Python y Flask. Las rutas del
módulo \textit{generator} se encuentran en
\path{app/features/generator/routes.py}. La lógica del asistente definida en
\path{app/features/generator/wizard.py} reconstruye y comprueba los valores
introducidos durante los pasos de configuración. La
Figura~\ref{fig:estructura_generator} muestra la estructura del módulo.

\begin{figure}[H]
    \centering
    \includegraphics[width=0.80\textwidth]{estructura_modulo_generator.png}
    \caption{Estructura del módulo generator de UVLHub. Elaboración propia.}
    \label{fig:estructura_generator}
\end{figure}

El estado del asistente se mantiene de forma temporal en la sesión de Flask.
Una vez consolidada la configuración, el endpoint
\texttt{/generator/random/params-json} la proporciona en formato JSON. En la
generación ordinaria, el navegador construye
\texttt{FmgeneratorModel} mediante \texttt{from\_flat\_dict} y ejecuta el
generador en Pyodide. Cuando se activa la opción de satisfacibilidad, el
navegador envía los parámetros al endpoint
\texttt{/generator/random/generate-sat}; el servidor genera los candidatos y
los comprueba mediante PySAT.

\subsection*{Generador compatible con Flamapy}

La lógica principal se ha implementado como un componente compatible con la
arquitectura de extensiones de Flamapy. El directorio \texttt{fm\_generator}
contiene \texttt{FmgeneratorModel} y las clases que representan los diferentes
bloques de configuración.

La operación \texttt{GenerateFeatureModel} recibe dicha configuración y
construye un objeto \texttt{FeatureModel} por ejecución. Su método
\texttt{execute} devuelve la propia operación y \texttt{get\_result} permite
recuperar el modelo obtenido. La transformación a UVL se realiza después
mediante \texttt{UVLWriter}. La
Figura~\ref{fig:estructura_fm_generator} muestra la estructura de la
herramienta.

\begin{figure}[H]
    \centering
    \includegraphics[width=0.80\textwidth]{estructura_fm_generator.png}
    \caption{Estructura de la herramienta generadora. Elaboración propia.}
    \label{fig:estructura_fm_generator}
\end{figure}

\subsection*{Capa de presentación}

La interfaz mantiene la arquitectura existente de UVLHub. Las vistas se
definen con plantillas Jinja2; JavaScript gestiona la interacción y Bootstrap
proporciona los componentes visuales. Los archivos de los seis pasos se
encuentran en \path{app/features/generator/templates/generator}, cuya
estructura se muestra en la Figura~\ref{fig:templates_generator}.

\begin{figure}[H]
    \centering
    \includegraphics[width=0.80\textwidth]{templates_generator.png}
    \caption{Plantillas HTML utilizadas para los pasos del asistente. Elaboración propia.}
    \label{fig:templates_generator}
\end{figure}

\subsection*{Ejecución mediante Pyodide y WebAssembly}

Pyodide~\cite{pyodide} permite ejecutar Python en el navegador mediante
WebAssembly~\cite{webassembly}. En el flujo ordinario, JavaScript prepara el
entorno, carga las dependencias necesarias y utiliza
\texttt{fmgen\_wrapper.py} para construir la configuración, generar cada
modelo y convertirlo a UVL. Los scripts y paquetes \textit{wheel} se
encuentran en \path{app/features/generator/assets/js}. La
Figura~\ref{fig:pyodide_assets} muestra su estructura sin desplegar las
dependencias empaquetadas ni los archivos del entorno WebAssembly.

\begin{figure}[H]
    \centering
    \includegraphics[width=0.80\textwidth]{pyodide_assets.png}
    \caption{Scripts utilizados para preparar y ejecutar el generador mediante Pyodide. Elaboración propia.}
    \label{fig:pyodide_assets}
\end{figure}

La generación con comprobación SAT sigue el recorrido del servidor: allí se
construye la configuración efectiva booleana, se generan los candidatos y se
comprueba su satisfacibilidad. Los modelos aceptados se devuelven al navegador
en formato JSON.

\subsection*{Infraestructura de desarrollo}

Durante el desarrollo local se ha utilizado Docker~\cite{docker} junto con
Docker Compose para ejecutar el entorno de UVLHub. Los archivos de
configuración se encuentran en el directorio \texttt{docker}. Para el
desarrollo se emplean \path{docker/docker-compose.dev.yml} y
\path{docker/images/Dockerfile.dev}; \path{docker/entrypoints} contiene, entre
otros, los scripts \path{development_entrypoint.sh},
\path{production_entrypoint.sh}, \path{webhook_entrypoint.sh} y
\path{worker_entrypoint.sh}.

También existen configuraciones de producción, como
\path{docker/docker-compose.prod.yml} y
\path{docker/docker-compose.prod.webhook.yml}.
Estos archivos pertenecen a UVLHub y su presencia no implica que todos esos
entornos se hayan desplegado durante el TFG. 

La Figura~\ref{fig:docker_structure} recoge la estructura del directorio.

\begin{figure}[H]
    \centering
    \includegraphics[width=0.80\textwidth]{docker_structure.png}
    \caption{Estructura del directorio \texttt{docker} de UVLHub. Elaboración propia.}
    \label{fig:docker_structure}
\end{figure}

Entre los servicios auxiliares se encuentra Redis, utilizado durante el
despliegue de revisiones para conservar temporalmente la sesión de Flask. La
generación ordinaria sigue ejecutándose en el navegador, mientras que la
comprobación SAT y sus reintentos se realizan en el servidor.

\subsection*{Suite de pruebas}

Las pruebas de UVLHub se ejecutan mediante \textit{Rosemary}, una interfaz de
línea de comandos incluida en el entorno de la aplicación~\cite{rosemary}.
La suite del módulo \textit{generator} incluye comprobaciones del asistente,
validación de parámetros, integración de la generación y pruebas de interfaz
mediante Selenium. Su estructura aparece en la
Figura~\ref{fig:tests_generator}.

\begin{figure}[H]
    \centering
    \includegraphics[width=0.80\textwidth]{tests_generator.png}
    \caption{Suite de pruebas del módulo generator. Elaboración propia.}
    \label{fig:tests_generator}
\end{figure}

La herramienta \textit{fm\_generator} dispone además de pruebas para su
modelo de configuración y su operación de generación, organizadas como
indica la Figura~\ref{fig:tests_fm_generator}.

\begin{figure}[H]
    \centering
    \includegraphics[width=0.80\textwidth]{tests_fm_generator.png}
    \caption{Suite de pruebas de \textit{fm\_generator}. Elaboración propia.}
    \label{fig:tests_fm_generator}
\end{figure}

\subsection*{Automatización CI/CD}

UVLHub dispone de flujos de GitHub Actions en su repositorio
\textit{upstream}. Sus comprobaciones pueden ejecutarse al presentar una
\textit{pull request} antes de una posible fusión. Estos flujos pertenecen a
UVLHub y no fueron creados durante este TFG. Su alcance y los comandos locales
utilizados para verificar el generador se describen en la
Sección~\ref{sec:ci}.

\section{Módulo de validación y lógica principal}
\label{sec:modulo_principal}

La validación se organiza en dos niveles. Durante el asistente se comprueban
los valores de cada paso y se mantiene el estado temporal en la sesión. Estas
tareas corresponden a \texttt{GeneratorWizardService}. Tras construir la
configuración completa, \texttt{FmgeneratorModel} comprueba las restricciones
de dominio.

La operación \texttt{GenerateFeatureModel} recibe esta configuración y
produce un \texttt{FeatureModel} por ejecución. Antes de generar el modelo,
\texttt{execute} vuelve a llamar a la validación del objeto:

\begin{lstlisting}[language=Python]
self.model.validate()
\end{lstlisting}

Esta comprobación protege la operación cuando se invoca directamente, tanto
desde el navegador como desde el servidor. En el flujo satisfacible, el
servidor construye una configuración efectiva limitada a construcciones
booleanas y comprueba mediante SAT cada candidato generado.

\subsection{Validación sintáctica y comprobaciones básicas}
\label{subsec:validacion_sintactica}

Las comprobaciones de los formularios se implementan principalmente en
\texttt{wizard.py}, mediante funciones asociadas a los seis pasos del
asistente:

\begin{lstlisting}[language=Python]
validate_step1_form()
validate_step2_form()
validate_step3_form()
validate_step4_form()
validate_step5_form()
validate_step6_form()
\end{lstlisting}

Estas funciones detectan valores ausentes o inválidos antes de continuar con
el asistente. También comprueban dependencias entre opciones y distribuciones
introducidas por el usuario. La validación de la configuración completa se
realiza posteriormente en \texttt{FmgeneratorModel}.

\subsubsection*{Validación de parámetros generales}

En el primer paso se comprueban el número de modelos y la semilla:

\[
1 \leq \text{\texttt{NUM\_MODELS}} \leq 1000,
\qquad
\text{\texttt{SEED}} \in \mathbb{Z}^{+}.
\]

\subsubsection*{Validación de rangos numéricos}

Los parámetros que representan cantidades positivas deben cumplir los
límites definidos para cada campo. En los pares de mínimo y máximo se
comprueba, de forma general:

\[
\text{\texttt{MIN\_VALUE}} \leq \text{\texttt{MAX\_VALUE}}.
\]

Esta relación se aplica, entre otros, al número de características,
restricciones y atributos. Algunos campos admiten el valor cero, por lo que
el límite inferior se comprueba según el significado de cada parámetro.

\subsubsection*{Validación de distribuciones probabilísticas}

Las probabilidades utilizadas en una decisión aleatoria deben pertenecer al
intervalo $[0,1]$ y formar las distribuciones definidas para cada grupo de
opciones:

\[
\forall p_i\in P,\quad 0\leq p_i\leq 1.
\]

En la interfaz, la distribución de relaciones jerárquicas admite una
diferencia de hasta $10^{-3}$ respecto a uno, como se especifica en RN-03.
La validación semántica de \texttt{FmgeneratorModel} utiliza una tolerancia
de $10^{-6}$ para comprobar las distribuciones normalizadas.

\subsubsection*{Validación de dependencias básicas}

Algunas opciones requieren que se activen antes los niveles de expresividad
correspondientes:

\[
\text{\texttt{FEATURE\_CARDINALITY}}
\Rightarrow \text{\texttt{ARITHMETIC\_LEVEL}}.
\]

\[
\text{\texttt{AGGREGATE\_FUNCTIONS}}
\Rightarrow \text{\texttt{ARITHMETIC\_LEVEL}}.
\]

\[
\text{\texttt{STRING\_CONSTRAINTS}}
\Rightarrow \text{\texttt{TYPE\_LEVEL}}.
\]

\subsubsection*{Validación de atributos manuales}

En el modo manual, cada atributo debe tener un nombre y una definición
compatible con su tipo. Si se indica una probabilidad de asociación, esta
debe estar comprendida entre cero y uno. Los dominios expresados mediante
un rango deben cumplir que el límite inferior no supere al superior.

La inclusión de atributos numéricos en restricciones requiere el nivel
aritmético. Para utilizar atributos de cadena en dichas restricciones deben
estar activos el nivel de tipos y las restricciones sobre cadenas.

\subsection{Validación semántica y restricciones de dominio}
\label{subsec:validacion_semantica}

La validación semántica se encapsula en \texttt{FmgeneratorModel}, definido
en \texttt{fm\_generator}. La Tabla~\ref{tab:generator_configs} resume los
bloques que agrupan su configuración.

\begin{table}[H]
    \centering
    \begin{tabular}{|c|c|}
        \hline
        \textbf{Configuración} & \textbf{Responsabilidad} \\
        \hline
        \texttt{NamingConfig} & Nombres y opciones de salida \\
        \hline
        \texttt{LevelConfig} & Niveles de expresividad UVL \\
        \hline
        \texttt{FeaturesConfig} & Parámetros de características \\
        \hline
        \texttt{HierarchyConfig} & Estructura jerárquica \\
        \hline
        \texttt{ConstraintsConfig} & Restricciones generadas \\
        \hline
        \texttt{AttributesConfig} & Gestión de atributos \\
        \hline
    \end{tabular}
    \caption{Bloques de configuración de \texttt{FmgeneratorModel}. Elaboración propia.}
    \label{tab:generator_configs}
\end{table}

El método \texttt{FmgeneratorModel.validate()} comprueba la coherencia global
antes de iniciar la generación. Entre otras condiciones, verifica las
dependencias entre niveles, los rangos numéricos y la normalización de las
distribuciones con una tolerancia de $10^{-6}$.

\subsubsection*{Coherencia de niveles UVL}

El nivel booleano constituye la base de los modelos generados. La activación
del nivel de tipos requiere el aritmético, y los niveles secundarios dependen
de sus niveles principales:

\[
\text{\texttt{BOOLEAN\_LEVEL}}=\text{\texttt{True}},
\]

\[
\text{\texttt{TYPE\_LEVEL}}
\Rightarrow\text{\texttt{ARITHMETIC\_LEVEL}}.
\]

\[
\text{\texttt{FEATURE\_CARDINALITY}}
\Rightarrow\text{\texttt{ARITHMETIC\_LEVEL}}.
\]

\[
\text{\texttt{AGGREGATE\_FUNCTIONS}}
\Rightarrow\text{\texttt{ARITHMETIC\_LEVEL}}.
\]

\[
\text{\texttt{STRING\_CONSTRAINTS}}
\Rightarrow\text{\texttt{TYPE\_LEVEL}}.
\]

\subsubsection*{Validación de la jerarquía de características}

La profundidad máxima se comprueba frente al número máximo de
características:

\[
1\leq\text{\texttt{MAX\_TREE\_DEPTH}}
\leq\text{\texttt{MAX\_FEATURES}}.
\]

Las probabilidades de selección de relaciones obligatorias, opcionales,
alternativas, disyuntivas y de cardinalidad de grupo deben formar una
distribución normalizada.

\subsubsection*{Validación de restricciones}

La configuración valida las probabilidades de los tipos de restricciones y
de los operadores disponibles. Una probabilidad positiva para restricciones
numéricas requiere el nivel aritmético; las restricciones sobre cadenas
requieren tanto el nivel de tipos como la opción correspondiente:

\[
P(\text{\texttt{Integer}})>0\ \lor\
P(\text{\texttt{Real}})>0
\Rightarrow\text{\texttt{ARITHMETIC\_LEVEL}},
\]

\[
P(\text{\texttt{String}})>0
\Rightarrow
\text{\texttt{TYPE\_LEVEL}}\land\text{\texttt{STRING\_CONSTRAINTS}}.
\]

Para las restricciones booleanas, las probabilidades de conjunción,
disyunción, implicación y equivalencia deben sumar uno. La probabilidad de
negación se valida de forma independiente dentro del intervalo $[0,1]$.

Las operaciones aritméticas se comprueban según las capacidades activadas.
Cuando se permiten funciones agregadas, sus probabilidades se incorporan a
la distribución utilizada para seleccionarlas. Los operadores de comparación
se validan de forma separada. De esta forma, solo se consideran opciones
compatibles con los niveles UVL habilitados.

\subsubsection*{Validación de atributos}

En el modo aleatorio se comprueba el rango del número de atributos y la
distribución de sus tipos:

\[
0\leq\text{\texttt{MIN\_ATTRIBUTES}}
\leq\text{\texttt{MAX\_ATTRIBUTES}}.
\]

En el modo manual se comprueba cada elemento de
\texttt{attributes\_list}, incluidos su dominio y, cuando se proporciona,
\texttt{attach\_probability}:

\[
0\leq
\text{\texttt{attach\_probability}}
\leq 1.
\]

La posibilidad de utilizar atributos en restricciones se valida de acuerdo
con los niveles aritmético y de tipos activos.

\subsubsection{Consumo por otros módulos}

La instancia validada de \texttt{FmgeneratorModel} se utiliza para la
generación. La operación \texttt{GenerateFeatureModel} construye la
jerarquía, asigna los atributos y genera las restricciones.
El archivo \path{tests/test_validators.py} incluye casos de rechazo de
configuraciones inconsistentes, como distribuciones no normalizadas,
dependencias incumplidas y rangos inválidos.

\section{Módulo de generación de artefactos}
\label{sec:generador}

La generación se encapsula en la operación
\texttt{GenerateFeatureModel}, perteneciente a \texttt{fm\_generator}.
Recibe una instancia de \texttt{FmgeneratorModel} y produce un objeto
\texttt{FeatureModel} en memoria:

\[
\text{\texttt{FmgeneratorModel}}
\xrightarrow{\text{\texttt{GenerateFeatureModel}}}
\text{\texttt{FeatureModel}}.
\]

El resultado contiene la jerarquía de características, sus relaciones,
atributos y restricciones en las estructuras del metamodelo de Flamapy. La
serialización a UVL y la asignación del nombre del archivo se realizan fuera
de la operación. Una invocación completa puede expresarse como:

\begin{lstlisting}[language=Python]
operation = GenerateFeatureModel().execute(
    config, index=index, attempt=attempt
)
feature_model = operation.get_result()
\end{lstlisting}

\subsection{Lógica principal de generación de modelos}
\label{subsec:calculos}

La lógica se implementa en \path{generate_feature_model.py}, dentro del
directorio de operaciones de \texttt{fm\_generator}. Los siguientes
fragmentos muestran la coordinación del proceso, la inicialización de la
aleatoriedad, la construcción de la jerarquía, la asignación de atributos y
la creación de restricciones.

\subsubsection*{Coordinación de la operación y semilla}

El método \texttt{execute} recibe la configuración y genera un único
\texttt{FeatureModel} por llamada. Primero valida los parámetros; después
construye la jerarquía, añade atributos y restricciones y guarda el
resultado. La operación devuelve \texttt{self} y el modelo se recupera
mediante \texttt{get\_result}. El Código~\ref{lst:execute-generador}
muestra esta secuencia.

\begin{lstlisting}[language=Python,
basicstyle=\ttfamily\footnotesize,breaklines=true,columns=fullflexible,
caption={Coordinación de la generación en \texttt{GenerateFeatureModel}.},
label={lst:execute-generador}]
def execute(
    self,
    model: FmgeneratorModel,
    index: int = 0,
    attempt: int = 0,
) -> Operation:

    self.model = model

    self.model.validate()
    self._seed_generation(index, attempt)

    fm, features = self._generate_hierarchy()
    self._assign_attributes(features)
    self._add_constraints(fm, features)

    setattr(fm, "uvl_includes", self._build_uvl_includes())

    self.result = fm
    return self

def get_result(self) -> FeatureModel | None:
    return self.result
\end{lstlisting}

El archivo define las constantes del Código~\ref{lst:constantes-generador}.
La primera separa las semillas de distintos intentos; la segunda fija la
probabilidad por defecto de incluir un atributo aleatorio entre los
candidatos de una restricción.

\begin{lstlisting}[language=Python,
basicstyle=\ttfamily\footnotesize,breaklines=true,columns=fullflexible,
caption={Constantes utilizadas para las semillas y los atributos aleatorios.},
label={lst:constantes-generador}]
SAT_SEED_STRIDE = 100000
RANDOM_ATTR_CONSTRAINT_PROB = 0.8
\end{lstlisting}

La semilla aplicada a un modelo del lote se calcula como:

\[
s_{\mathrm{efectiva}}=
s_{\mathrm{base}}+i+100000a,
\]

donde $i$ representa el índice del modelo y $a$ el número de intento.
El Código~\ref{lst:semilla-generador} muestra su implementación.

\begin{lstlisting}[language=Python,
basicstyle=\ttfamily\footnotesize,breaklines=true,columns=fullflexible,
caption={Cálculo de la semilla para cada modelo e intento.},
label={lst:semilla-generador}]
def _seed_generation(self, index: int, attempt: int) -> None:
    seed_value = self.model.seed + index + (attempt * SAT_SEED_STRIDE)
    random.seed(seed_value)
\end{lstlisting}

La misma configuración, semilla, versión del generador y condiciones de
ejecución permiten repetir una ejecución. Los reintentos con distintos
valores de \texttt{attempt} emplean secuencias pseudoaleatorias diferentes.

\subsubsection*{Construcción de la jerarquía}

La raíz se denomina \texttt{F0}. El método
\texttt{\_generate\_hierarchy} construye las demás características dentro
del rango configurado, las distribuye entre los niveles permitidos por
\texttt{max\_tree\_depth} y crea sus relaciones con los padres. El
Código~\ref{lst:relaciones-generador} muestra cómo selecciona los tipos de
relación y conecta los elementos de un nivel.

\begin{lstlisting}[language=Python,
basicstyle=\ttfamily\footnotesize,breaklines=true,columns=fullflexible,
caption={Selección y conexión de relaciones entre características.},
label={lst:relaciones-generador}]
def _select_relation_types(self, total: int) -> list[str]:
    return random.choices(
        population=["mand", "opt", "alt", "or", "group"],
        weights=[
            self.model.hierarchy.dist_mandatory,
            self.model.hierarchy.dist_optional,
            self.model.hierarchy.dist_alternative,
            self.model.hierarchy.dist_or,
            self.model.hierarchy.dist_group_cardinality,
        ],
        k=total,
    )

def _add_relations_to_level(
    self,
    parents: list[Feature],
    children: list[Feature],
) -> None:
    total = len(children)
    relation_types = self._select_relation_types(total)
    random.shuffle(relation_types)

    pool = children[:]
    parent_index = 0

    while pool:
        relation_kind = relation_types[parent_index % len(relation_types)]
        parent = parents[parent_index % len(parents)]
        parent_index += 1

        size = self._determine_group_size(len(pool))
        group = [pool.pop() for _ in range(size)]

        relations = self._create_relation(parent, group, relation_kind)

        for relation in relations:
            parent.add_relation(relation)

            for child in relation.children:
                child.parent = parent
                self._maybe_apply_feature_cardinality(child)
\end{lstlisting}

\texttt{random.choices} obtiene tipos de relación según los pesos indicados
por el usuario. Después se mezclan y se agrupan los hijos. Por ello, las
probabilidades influyen en la selección, pero no garantizan un número exacto
de relaciones de cada tipo. \texttt{\_create\_relation} establece los
límites de cada relación: obligatoria $[1,1]$, opcional $[0,1]$,
alternativa $[1,1]$ para el grupo y disyuntiva $[1,n]$ para un grupo de
$n$ hijos. En la cardinalidad de grupo, los límites se eligen de acuerdo
con el tamaño del grupo.

El método \texttt{\_maybe\_apply\_feature\_type} asigna el tipo a cada
característica. Si el nivel de tipos está
desactivado, las características son booleanas. Cuando se activa, el método
evalúa los pesos configurados para los tipos booleano, entero, real y cadena,
y elige entre los candidatos obtenidos. Si no resulta seleccionado ningún
tipo alternativo, se conserva el booleano.

\subsubsection*{Asignación de atributos}

\texttt{\_assign\_attributes} distingue entre generación aleatoria y
definición manual. En el primer caso se decide el número de atributos y se
reparten entre las características; en el segundo se recorren las
definiciones introducidas por el usuario y se aplica su probabilidad de
asociación. El Código~\ref{lst:atributos-generador} muestra la selección del
modo.

\begin{lstlisting}[language=Python,
basicstyle=\ttfamily\footnotesize,breaklines=true,columns=fullflexible,
caption={Selección del modo de asignación de atributos.},
label={lst:atributos-generador}]
def _assign_attributes(self, features: list[Feature]) -> None:
    if self.model.attributes.random_attributes:
        self._generate_random_attributes(features)
    else:
        self._assign_manual_attributes(features)
\end{lstlisting}

En el modo aleatorio, cuando el nivel aritmético está activado y la
distribución permite crear tipos numéricos, se reserva una cantidad inicial
de atributos numéricos. Si $v_{\min}$ representa
\texttt{min\_vars\_per\_constraint}, $r$ representa
\texttt{extra\_constraint\_representativeness}, $A$ el número de atributos
solicitados y $F$ las características disponibles, el cálculo utilizado es:

\[
A_{\mathrm{num}}=
\min\left(
\left\lceil\frac{v_{\min}}{r}\right\rceil,\ A,\ F
\right).
\]

Esta reserva proporciona candidatos para las restricciones numéricas. Los
restantes tipos de atributo se eligen con \texttt{random.choices} según la
distribución configurada. Sus dominios y valores se materializan mediante
objetos \texttt{Domain}, \texttt{Range} y \texttt{Attribute} antes de
asociarse a las características.

\subsubsection*{Construcción de restricciones}

\texttt{\_add\_constraints} organiza las características y los atributos
en conjuntos de candidatos booleanos, numéricos y de cadenas. En modo
\texttt{ensure\_satisfiable},
\texttt{\_constraints\_must\_be\_boolean\_only} limita los candidatos a
expresiones booleanas. Esta limitación permite aplicar el comprobador SAT,
pero la satisfacibilidad se verifica posteriormente en el servidor.

\texttt{extra\_constraint\_representativeness} limita cuántas veces
puede aparecer una característica en una restricción. El
Código~\ref{lst:muestreo-restricciones} muestra cómo se seleccionan
variables respetando este máximo y el límite de características distintas
que recibe el método.

\begin{lstlisting}[language=Python,
basicstyle=\ttfamily\footnotesize,breaklines=true,columns=fullflexible,
caption={Muestreo de variables con límite de repeticiones.},
label={lst:muestreo-restricciones}]
def _sample_keys_with_ecr(
    self,
    groups: dict[str, list[str]],
    target_occurrences: int,
    max_repetitions: int,
    max_features_param: int,
    feature_usage: dict[str, int] | None = None,
    selected_features: set[str] | None = None,
) -> list[str] | None:
    if target_occurrences < 1:
        return None

    feature_usage = {} if feature_usage is None else feature_usage
    selected_features = set() if selected_features is None else selected_features

    selected_keys: list[str] = []

    for _ in range(target_occurrences):
        allowed_feature_ids = [
            feature_id
            for feature_id in groups
            if feature_usage.get(feature_id, 0) < max_repetitions
            and (
                feature_id in selected_features
                or len(selected_features) < max_features_param
            )
        ]

        if not allowed_feature_ids:
            return None

        feature_id = random.choice(allowed_feature_ids)
        selected_keys.append(random.choice(groups[feature_id]))
        feature_usage[feature_id] = feature_usage.get(feature_id, 0) + 1
        selected_features.add(feature_id)

    random.shuffle(selected_keys)
    return selected_keys
\end{lstlisting}

Si $g$ es el número de características candidatas, $f_{\max}$ el máximo
de características distintas y $r$ el máximo de repeticiones por
característica, la capacidad calculada para ese conjunto es:

\[
C_{\max}=\min(g,f_{\max})\,r.
\]

Cuando no hay suficientes candidatos para alcanzar el mínimo de variables
exigido, se omite la restricción. En caso contrario, se construyen los
predicados mediante nodos \texttt{Node} y operaciones
\texttt{ASTOperation} del árbol de sintaxis abstracta. El
Código~\ref{lst:append-restriccion} muestra cómo se incorpora una
restricción construida al modelo.

\begin{lstlisting}[language=Python,
basicstyle=\ttfamily\footnotesize,breaklines=true,columns=fullflexible,
caption={Incorporación de una restricción al modelo de características.},
label={lst:append-restriccion}]
if root is None:
    continue

fm.ctcs.append(Constraint(name=f"ctc{index}", ast=AST(root)))
\end{lstlisting}

\texttt{min\_constraints} y \texttt{max\_constraints} determinan el
número de intentos de creación. El resultado puede contener menos
restricciones efectivas cuando no existen candidatos compatibles o no se
consigue construir un predicado válido.

\subsubsection*{Directivas de expresividad UVL}

\texttt{\_build\_uvl\_includes} reúne las directivas de los niveles
secundarios activados. Si se habilitan a la vez la cardinalidad de
característica y las funciones agregadas, se agrupan en
\texttt{Arithmetic.*}, como recoge el
Código~\ref{lst:includes-generador}.

\begin{lstlisting}[language=Python,
basicstyle=\ttfamily\footnotesize,breaklines=true,columns=fullflexible,
caption={Construcción de las directivas \texttt{include} del resultado.},
label={lst:includes-generador}]
def _build_uvl_includes(self) -> list[str]:
    includes: list[str] = []

    if self.model.levels.group_cardinality:
        includes.append("Boolean.group-cardinality")

    feature_cardinality = self.model.levels.feature_cardinality
    aggregate_functions = self.model.levels.aggregate_functions

    if feature_cardinality and aggregate_functions:
        includes.append("Arithmetic.*")
    else:
        if aggregate_functions:
            includes.append("Arithmetic.aggregate-function")
        if feature_cardinality:
            includes.append("Arithmetic.feature-cardinality")

    if self.model.levels.type_level and self.model.levels.string_constraints:
        includes.append("Type.string-constraints")

    return includes
\end{lstlisting}

Estas directivas se guardan en el \texttt{FeatureModel} y se utilizan al
transformarlo posteriormente a texto UVL mediante \texttt{UVLWriter}.

\subsection{Síntesis del artefacto final generado}
\label{subsec:sintesis}

\subsubsection*{Preparación de la configuración}

Al finalizar el asistente, el navegador obtiene los parámetros en formato
JSON mediante \texttt{/generator/random/params-json}. Si el usuario solicita
la generación ordinaria, \texttt{fmgen\_wrapper.py} construye
\texttt{FmgeneratorModel} dentro de Pyodide mediante
\texttt{from\_flat\_dict}. Si solicita la comprobación de
satisfacibilidad, envía los parámetros al servidor, donde se construye una
configuración efectiva limitada al nivel booleano. Se desactivan los niveles
aritmético y de tipos, sus niveles secundarios, la cardinalidad de grupo y
la generación de atributos.

\subsubsection*{Generación del modelo intermedio}

\texttt{GenerateFeatureModel} construye un \texttt{FeatureModel} por
ejecución. En el navegador, cada resultado obtenido se prepara para su
serialización. En el servidor, cada resultado es un candidato que se
comprueba mediante PySAT. Si no es satisfacible, se genera otro candidato
con un número de intento distinto, hasta un máximo de 20 intentos por
modelo. Si se agotan los intentos, se produce
\texttt{SatisfiableModelGenerationError} y no se devuelve el último
candidato como resultado válido.

\subsubsection*{Transformación a formato UVL}

\texttt{UVLWriter} transforma los objetos \texttt{FeatureModel} aceptados
en texto UVL. La operación \texttt{GenerateFeatureModel} termina al producir
el modelo en memoria; no realiza esta transformación.

\subsubsection*{Entrega del resultado}

En el flujo ordinario, el navegador dispone directamente de los textos UVL
generados. En el flujo satisfacible, el servidor devuelve los resultados
aceptados en formato JSON. La interfaz permite visualizarlos y descargarlos,
individualmente o agrupados en un ZIP según la opción disponible. Esta
funcionalidad no añade almacenamiento permanente de los modelos.

\section{Rutas HTTP y API del generador}
\label{sec:api}

\subsection{Endpoints implementados}
\label{subsec:endpoints}

El módulo ofrece rutas para mostrar y procesar los formularios del
asistente, además de endpoints que intercambian información en JSON.

\subsubsection*{Endpoints de entrada al generador}

\begin{itemize}
    \item \textbf{\texttt{GET /generator}:} inicia el asistente y muestra
    la página inicial.

    \item \textbf{\texttt{GET /generator/random}:} dirige al primer paso de
    la configuración para la generación aleatoria.
\end{itemize}

\subsubsection*{Endpoints del asistente de configuración}

Cada paso admite peticiones \texttt{GET} para mostrar su formulario y
\texttt{POST} para procesar los valores introducidos:

\begin{itemize}
    \item \textbf{\texttt{GET/POST /generator/random/step1}:} número de
    modelos, semilla y opciones generales.

    \item \textbf{\texttt{GET/POST /generator/random/step2}:} niveles de
    expresividad UVL.

    \item \textbf{\texttt{GET/POST /generator/random/step3}:} estructura de
    características y relaciones jerárquicas.

    \item \textbf{\texttt{GET/POST /generator/random/step4}:} restricciones
    y operadores.

    \item \textbf{\texttt{GET/POST /generator/random/step5}:} atributos.

    \item \textbf{\texttt{GET/POST /generator/random/step6}:} opciones
    finales y presentación de la salida. La ejecución del generador se
    inicia posteriormente desde el navegador o mediante la ruta SAT.
\end{itemize}

\subsubsection*{Endpoints auxiliares y de generación}

\begin{itemize}
    \item \textbf{\texttt{GET /generator/random/params-json}:} proporciona
    al navegador la configuración consolidada del asistente en formato JSON.

    \item \textbf{\texttt{POST /generator/random/summary-refresh/<step>}:}
    actualiza el resumen de la configuración mostrado en el asistente.

    \item \textbf{\texttt{POST /generator/random/generate-sat}:} recibe
    los parámetros JSON, construye la configuración efectiva booleana y
    genera candidatos en el servidor. Comprueba cada candidato mediante PySAT,
    realiza hasta 20 intentos por modelo y devuelve en formato JSON los modelos
    satisfacibles. Si se agotan los intentos, devuelve el error correspondiente.
\end{itemize}

La generación ordinaria se ejecuta en Pyodide después de consultar
\texttt{params-json}; no requiere una petición al servidor para generar
cada \texttt{FeatureModel}.

\subsection{Descarga de modelos UVL generados}
\label{subsec:descarga}

Una vez obtenidos los textos UVL, la interfaz permite consultar los modelos
y descargarlos, incluidos los lotes ofrecidos en ZIP. El servidor no conserva
permanentemente los modelos generados. Para repetir la generación se utilizan
la misma configuración, semilla, versión del generador y condiciones de
ejecución.

\section{Integración continua}
\label{sec:ci}

El repositorio \textit{upstream} de UVLHub utiliza GitHub Actions. Cuando se
propone integrar cambios mediante una \textit{pull request}, las
comprobaciones configuradas allí sirven como control previo a una posible
fusión. Los flujos de la Figura~\ref{fig:workflows-upstream} pertenecen a
UVLHub y no fueron desarrollados durante este TFG. La captura muestra sus
nombres, pero no permite identificar sus disparadores ni los comandos de
cada archivo \texttt{.yml}.

\begin{figure}[H]
    \centering
    \includegraphics[width=0.50\textwidth]{workflows.png}
    \caption{\textit{Workflows} visibles en GitHub Actions del repositorio \textit{upstream} de UVLHub. Captura facilitada por el autor.}
    \label{fig:workflows-upstream}
\end{figure}

En la captura aparecen \textit{Commits Syntax Checker}, \textit{Pytest} y
\textit{Python Lint}, relacionados por sus nombres con comprobaciones
previas a la integración. También se muestran \textit{Deploy on Webhook},
\textit{Publish image in Docker Hub}, \textit{Dependabot Updates} y
\textit{Dependency Graph}. Su presencia no demuestra que todos se ejecuten
al abrir una \textit{pull request} ni que las pruebas desplieguen
automáticamente la versión desarrollada en este TFG.

Las comprobaciones de este trabajo se ejecutaron localmente. Desde la raíz
del repositorio, los siguientes comandos permiten levantar el entorno de
desarrollo, ejecutar la suite del módulo integrado, medir su cobertura y
ejecutar las pruebas del paquete independiente:

\begin{verbatim}
docker compose -f docker/docker-compose.dev.yml up -d --build

docker compose -f docker/docker-compose.dev.yml exec web \
  rosemary test generator --all

docker compose -f docker/docker-compose.dev.yml exec web \
  rosemary coverage generator --all

pytest fm_generator/tests -v \
  --cov=fm_generator \
  --cov-report=term-missing \
  --cov-report=html
\end{verbatim}

El comando de cobertura de UVLHub también ejecuta sus pruebas, por lo que
sus resultados no se suman como casos nuevos a los de
\texttt{rosemary test}. Los resultados y el desglose de cobertura del
código de producción se presentan en la Sección~\ref{sec:cobertura}.
El despliegue en Render descrito en el
Apéndice~\ref{app:manual_instalacion} se utilizó para las revisiones de los
tutores y no se considera una ejecución de los flujos de
\textit{upstream}.

La implementación descrita comprende la configuración del asistente, la
generación ordinaria en Pyodide y la generación satisfacible en el servidor.
El siguiente capítulo presenta las pruebas utilizadas para verificar el
comportamiento de estos componentes.
